\section{Einleitung und Problemhintergrund}

Statistik Austria meldete 2025 einen erneuten Anstieg des Pkw-Bestands um 54.208 Fahrzeuge \citep{statistikaustriaPkwBestand2025Um2026} sowie die höchste Zahl an Neuzulassungen seit 2019 \citep{statistikaustria2025VielePkwNeuzulassungen2026}.
Diese Entwicklung verdeutlicht den anhaltenden Trend steigender Motorisierung in Österreich.
Mit mehr als 7,5 Millionen zugelassenen \acp{kfz} wurde zwischen 1990 und 2025 ein Höchststand erreicht, wobei besonders große Fahrzeuge wie Geländewagen und \acp{suv} zunehmend öffentlichen Raum beanspruchen \citep{hatzlDatenZeigenImmer2025}.
Dieser Trend erhöht den Druck auf städtische Infrastrukturen und macht eine effizientere Nutzung vorhandener Verkehrsflächen erforderlich \citep{jaschinskyWarumOsterreichsStadte2025}.

Ein zentraler Aspekt dieser Entwicklung ist die zunehmende Schwierigkeit, geeignete Parkmöglichkeiten zu finden.
Während Wien selbst einen vergleichsweise geringeren \acs{kfz}-Bestand aufweist, pendeln Erwerbstätige regelmäßig aus dem Umland in die Stadt \citep{statistikaustriaKfzBestand2026}.
Aufgrund teilweise unzureichender Anbindungen an den öffentlichen Verkehr erfolgt diese Mobilität häufig mit dem privaten \acs{kfz} \citep{aknoAKNiederosterreichPendleranalyse20252025}.
Die Suche nach einem Parkplatz endet dabei häufig in Parkhäusern, die nur teilweise mit Parkleitsystemen ausgestattet sind und meist nur statische oder lokal begrenzte Informationen liefern \citep{stadtwienParkleitsystemFuerWien2026}. Für Anreisende bedeutet dies häufig wiederholtes Anfahren mehrerer Parkhäuser sowie zusätzlichen Zeitverlust bei der Suche nach einem freien Stellplatz.
Der dadurch entstehende Suchverkehr verstärkt zugleich die Belastung urbaner Infrastrukturen und führt zu vermeidbaren Emissionen sowie höherer Lärmbelastung in den betroffenen Stadtteilen.

Auf der Angebotsseite stehen Parkhausbetreiber unter wachsendem Druck, ihre Stellplätze effizient auszulasten und sich gegenüber alternativen Mobilitätsangeboten zu behaupten \citep{wangEconomicsSmartParking2025}. \acp{sps} versprechen genau diese Effizienzgewinne durch Echtzeit-Belegungsdaten \citep{alaanzyRealTimeSmart2025}, dynamische Steuerung und eine bessere Auffindbarkeit freier Stellplätze \citep{sarkerSmartParkingSystem2020}. Gleichzeitig erfordern solche Systeme jedoch Investitionen in Sensorik, Cloud-Infrastruktur und laufenden Betrieb, deren wirtschaftliche Rentabilität für viele Betreiber unklar bleibt.

Hinzu kommt ein gesellschaftlicher Aspekt: das aktive Suchen nach Parkinformationen während der Fahrt erhöht das Ablenkungspotential mobiler Endgeräte und beeinträchtigt nachweislich Fahrverhalten und Aufmerksamkeit \citep{oviedo-trespalaciosUnderstandingImpactsMobile2016}.
In Österreich ist rund ein Drittel der Verkehrsunfälle mit Todesfolge auf Unachtsamkeit und Ablenkung zurückzuführen \citep{bmi397VerkehrstoteAuf2026}.
Ein verlässliches Smart-Parking-System kann hier nicht nur den Betreibern, sondern auch den Endnutzern einen Mehrwert in Form höherer Verkehrssicherheit bieten.

Im Rahmen dieser Arbeit wird daher ein prototypisches Smart-Parking-System entworfen, das auf einer \ac{aws}-basierten Infrastruktur mit \ac{iot} Core, \ac{mqtt} und serverlosen Lambda-Funktionen aufbaut. Zur Veranschaulichung wird ein Prototyp entwickelt: ein 3D-gedrucktes Modell einer Parkgarage mit Matchbox-Fahrzeugen sowie eine virtuelle Parkgarage mit Ampelsystem. Auf dieser Basis wird eine Kosten-Nutzen-Analyse aus Sicht eines Parkhausbetreibers entwickelt, die Investitions- und Betriebskosten dem erwarteten wirtschaftlichen Nutzen gegenüberstellt.

Der weitere Aufbau dieser Arbeit gliedert sich wie folgt: Kapitel \ref{sec:related_work} fasst den Stand des Wissens zu \acp{sps}, \ac{iot}-basierten Architekturen sowie zu wirtschaftlichen Bewertungen vergleichbarer Systeme zusammen.
Zudem wird der eigene Beitrag dieser Arbeit von bestehenden Forschungsergebnissen abgegrenzt.
Anschließend stellt Kapitel \ref{sec:methodology} den methodischen Ansatz von Drive \& Decide vor.
Dort werden der zugrunde liegende Use Case, der technische Aufbau des Prototypen mit \ac{aws}-Backend sowie das Kosten-Nutzen-Modell aus Sicht eines Parkhausbetreibers im Detail beschrieben. Außerdem wird erläutert, anhand welcher Kennzahlen die wirtschaftliche Sinnhaftigkeit der Lösung bewertet werden soll.
Kapitel \ref{sec:conclusion} fasst abschließend die zentralen Erkenntnisse zusammen, leitet die wichtigsten Implikationen für Parkhausbetreiber ab und gibt einen Ausblick auf weiterführende Arbeiten.
